% !TeX root = ../main.tex
\section{Type System and Soundness Theorem}
\label{sec:type-system}

% !TeX root = ../main.tex
\begin{figure}[H]
\centering
\small
\begin{mathpar}
 \inferrule*[right=Refl]{ }{\tau<:\tau}
 \and
 \inferrule*[right=Trans]{\tau<:\sigma\\\sigma<:\rho}{\tau<:\rho}
 \and
 \inferrule*[right=Real-Sub]{ }{\tyR\G<:\tyR\E}
 \and
 \inferrule*[right=Arrow]{\sigma_1<:\tau_1\\\tau_2<:\sigma_2}
   {\tau_1\to\tau_2<:\sigma_1\to\sigma_2}
 \and
 \inferrule*[right=Product]{\tau_1<:\sigma_1\\\tau_2<:\sigma_2}
   {\tau_1\times\tau_2<:\sigma_1\times\sigma_2}
 \and
 \inferrule*[right=Sum]{\tau_1<:\sigma_1\\\tau_2<:\sigma_2}
   {\tau_1+\tau_2<:\sigma_1+\sigma_2}
 \and
 \inferrule*[right=List]{\tau<:\sigma}{\tyL\tau<:\tyL\sigma}
\end{mathpar}
\begin{mathpar}
 \inferrule*[right=Sub]{\typing e\tau\\\tau<:\sigma}{\typing e\sigma}
 \and
 \inferrule*[right=Real]{r\in\Real}{\typing r{\tyR m}}
 \and
 \inferrule*[right=Reject]{ }{\typing{\kw{reject}}\tau}
 \and
 \inferrule*[right=Neg]{\typing e{\tyR m}}{\typing{-e}{\tyR m}}
 \and
 \inferrule*[right=Add]{\typing{e_1}{\tyR m}\\\typing{e_2}{\tyR m}}
   {\typing{e_1+e_2}{\tyR m}}
 \and
 \inferrule*[right=Mul]{\typing{e_1}{\tyR\G}\\\typing{e_2}{\tyR m}}
   {\typing{e_1\cdot e_2}{\tyR m}}
 \and
 \inferrule*[right=Div]{\typing{e_1}{\tyR m}\\\typing{e_2}{\tyR\G}}
   {\typing{e_1/e_2}{\tyR m}}
 \and
 \inferrule*[right=Cmp]{\typing{e_1}{\tyR\G}\\\typing{e_2}{\tyR\G}}
   {\typing{e_1<e_2}{\tyB}}
\end{mathpar}
\begin{mathpar}
 \inferrule*[right=Uniform]{\typing a{\tyR m}\\\typing b{\tyR m}}
   {\typing{\draw{uniform}\alpha{a,b}}{\tyR m}}
 \and
 \inferrule*[right=Gaussian]{\typing u{\tyR m}\\\typing v{\tyR\G}}
   {\typing{\draw{gaussian}\alpha{u,v}}{\tyR m}}
 \and
 \inferrule*[right=Poisson]{\typing\lambda{\tyR m}}
   {\typing{\draw{poisson}\alpha\lambda}{\tyR m}}
 \and
 \inferrule*[right=Exponential]{\typing\lambda{\tyR\G}}
   {\typing{\draw{exponential}\alpha\lambda}{\tyR m}}
 \and
 \inferrule*[right=Bernoulli]{\typing p{\tyR m}}
   {\typing{\draw{bernoulli}\alpha p}{\tyR m}}
 \and
 \inferrule*[right=Discrete]{\typing p{\tyL{\tyR m}}}
   {\typing{\draw{discrete}\alpha p}{\tyR m}}
 \and
 \inferrule*[right=Beta]{\typing a{\tyR\G}\\\typing b{\tyR\G}}
   {\typing{\draw{beta}\alpha{a,b}}{\tyR m}}
 \and
 \inferrule*[right=Gamma]{\typing k{\tyR m}\\\typing\lambda{\tyR\G}}
   {\typing{\draw{gamma}\alpha{k,\lambda}}{\tyR m}}
\end{mathpar}
\caption{Subtyping and typing rules, with $m\in\{\G,\E\}$ and
$\alpha\in\{m,\M\}$. Rules for variables, unit and Boolean literals, functions,
recursion, application, let bindings, conditionals, pairs, projections, sum
injections and case analysis, and list constructors and case analysis are omitted
because they are standard.}
\label{fig:typing}
\Description{Subtyping permits G reals where E reals are expected. Multiplication requires a G left operand, division a G denominator, and comparison two G operands. Distribution rules share the same premises for sampling and mean evaluation.}
\end{figure}


We use the type system in \cref{fig:typing} to check that mode $\E$ sampling is safe.
The type system checks that $\E$ samples are used ``expectation-affinely'':
they influence the expectation value of the final returned value in an affine way.
For example:
\begin{itemize}
    \item The arguments of addition are at $\E$-mode when the sum itself is.
    \item Multiplication requires one $\G$-operand and allows one $\E$-operand.
    \item Division requires a $\G$ denominator but allows an $\E$ numerator.
    \item The distribution rules allow $\E$ parameters where the mean is
affine in those parameters.
\end{itemize}

The aim of the type system is a soundness theorem for the determinization
transformation, which replaces $\E$-mode sampling with $\M$-mode deterministic mean computation.

\begin{theorem}[Expectation preservation]
\label{thm:expectation}
Suppose $\vdash e:\tyR\E$ and $e$ is domain-safe. If $\Expect[e]$ is defined,
then $\Expect[\trans e]$ is defined and
\[
 \Expect[\trans e]=\Expect[e].
\]
\end{theorem}
where
\begin{itemize}
    \item we write
$
 \Expect[e]\mathrel{:=}\int_{\Real} r\,\law e(dr)
$
for the unnormalized expected output of $e$.
    \item a program is \emph{domain-safe} if,
almost surely, every executed distribution has parameters in its domain and
every executed division has a nonzero divisor.
\end{itemize}

We also show that determinization can only improve the variance:

\begin{corollary}[Variance non-increase]
\label{cor:variance}
Suppose $\vdash e:\tyR\E$ and $e$ is domain-safe. If
$\law e(\Real)>0$ and $\int_{\Real}r^2\,\law e(dr)<\infty$, then $\trans e$
also has positive probability of returning and a finite second moment, and
\[
 \Var_{\mathrm{ret}}[\trans e]\le\Var_{\mathrm{ret}}[e].
\]
\end{corollary}
where
\begin{itemize}
    \item $\Var_{\mathrm{ret}}$ is the variance conditioned on the program succesfully returning (i.e., discarding rejected trials and infinite loops)
\end{itemize}

The reader may note that both theorems require the expectation or variance of the
input program to be well defined (which makes sense for theorems that state that expectation is preserved and variance only decreases).
We give the mathematical background below.
\Cref{sec:traces} gives refined theorems that apply even when the input program's
expectation is not well defined, and state that the expectation, although undefined, is still preserved in a finer grained sense.

\paragraph{When is an expectation defined?}
For a nonnegative quantity, the integral is always defined (by taking a supremum), although it may be
infinite. To integrate a signed output, separate its positive and negative parts:
\[
 r^+=\max(r,0),\qquad r^-=\max(-r,0),\qquad
 A_e=\int_{\Real}r^+\,\law e(dr),\quad
 B_e=\int_{\Real}r^-\,\law e(dr).
\]
Then $\Expect[e]=A_e-B_e$. If both parts are finite, the expectation is finite.
If exactly one part is infinite, the expectation is $+\infty$ or $-\infty$.
If both are infinite, the expectation is undefined: the two infinities cannot
be cancelled. In particular, having a finite expectation is equivalent to
\[
 \int_{\Real}|r|\,\law e(dr)=A_e+B_e<\infty.
\]
The absolute-value integral is nonnegative, so the question is whether it is
\emph{finite}, not whether it is defined.

For example, consider
\begin{verbatim}
let u = uniform[G](0, 1) in
let s = 2 * bernoulli[G](1/2) - 1 in
s / u + uniform[E](-1, 1)
\end{verbatim}
The program returns almost surely and is domain-safe: the exceptional event
$u=0$ has probability zero. Yet $\int_0^1 du/u=\infty$, and the random sign makes
both the positive and negative parts infinite, so the program's expectation is undefined. In contrast, $1/u$
alone has the defined expectation $+\infty$.

Determinization still removes the centered $\E$ noise and returns $s/u$. For almost
every fixed $\G$ trace, $u$ and $s$ are fixed, the source has the finite
conditional mean $s/u$, and the target returns exactly that number. This is the
property formalized in \cref{sec:traces}, which applies to all programs, regardless
of whether their expectation is defined.

\paragraph{Domain safety and missing probability.}
Our theorems require that the program uses the distributions in a domain-safe way.
For example, the following program is not domain safe:
\begin{verbatim}
let u = gaussian(0, 1) in uniform(0, u)
\end{verbatim}
because $u$ may be negative, and then the uniform distribution's parameters are invalid.

\paragraph{Rejection and divergence}
Rejection (``observe'' in other languages) is permitted:
we formally model $\kw{reject}$ as an infinite loop that contributes no output.

Let $d_e$ denote the probability of divergence, including rejection. It is the
limit, as $n$ grows, of the probability of executing $n$ steps and still not
being a value.

The following theorem states the analogue of conventional type safety,
and safety of the determinization transformation.

\begin{proposition}[Probability preservation]
\label{prop:probability}
Suppose $\vdash e:\tyR\E$ and $e$ is domain-safe. Then $\trans e$ is domain-safe,
and
\[
 \law e(\Real)+d_e=1,
 \qquad
 \law{\trans e}(\Real)=\law e(\Real).
\]
\end{proposition}

Thus the missing output mass accounts precisely for divergence and rejection;
neither the original program nor the transformed program loses probability mass due to getting stuck.

When $q_e=\law e(\Real)>0$, the distribution of successfully returned outputs
is $\law e/q_e$. Its expectation and variance are
\[
 \begin{aligned}
  \Expect_{\mathrm{ret}}[e]&=\frac{\Expect[e]}{q_e},\\
  \Var_{\mathrm{ret}}[e]&=\frac1{q_e}\int_{\Real}
    \bigl(r-\Expect_{\mathrm{ret}}[e]\bigr)^2\,\law e(dr).
 \end{aligned}
\]
The expectation is used when defined, and the variance when the second moment
is finite. Since determinization preserves $q_e$, \cref{thm:expectation} also
preserves $\Expect_{\mathrm{ret}}[e]$.
